% !TeX root = ../main.tex
\begin{cabstract}
粒子滤波是一种使用经验测度来拟合后验测度的方法，由于其无需正态分布假设，因而广泛适用于非线性系统的状态估计。然而，在实际应用中，传统粒子滤波存在一些固有局限：弱过程噪声乃至零过程噪声系统下，重采样步骤会造成单点塌缩；高维观测场景下，似然权重会集中于个别样本，也就是粒子贫化。

本文引入最优输运与神经网络方法来解决上述局限。在重采样步骤，使用输运变换来代替随机重采样，是解决退化问题的经典方法，但本文额外研究了误差的正交分解理论，具有可解释性。对于高维观测场景，本文在经典的条件输运对抗学习方法基础上，构建了摊销条件输运粒子滤波算法，进而同时提升了精度和在线速度。实验选取了三类问题：一维非线性方程、一维零过程噪声混沌系统及含噪手写数字逐窗识别。实验表明，输运方法为粒子滤波增强了稳健性，而离线摊销大幅减少了神经网络的在线推理开销。
\end{cabstract}
\begin{eabstract}
Particle filtering uses an empirical measure to fit the posterior measure; since it requires no normality assumption, it applies broadly to state estimation in nonlinear systems. In practice, however, traditional particle filters have inherent limitations: under weak or even zero process noise, the resampling step causes particles to collapse to a single point; in high-dimensional observation settings, likelihood weights concentrate on just a few samples, i.e., particle impoverishment.

This paper introduces optimal transport and neural network methods to address these limitations. Replacing random resampling with a transport transformation is a classic approach to handling degeneracy, but here we additionally develop an orthogonal error decomposition theory for clear interpretability. For high-dimensional observations, building on classical conditional transport adversarial learning, we construct an amortized conditional transport particle filter that improves both accuracy and online speed. Experiments cover three problems: a one-dimensional nonlinear equation, a one-dimensional chaotic system with zero process noise, and window-by-window recognition of noisy handwritten digits. The results show that transport methods improve the robustness of particle filtering, while offline amortization substantially reduces the online inference overhead of the neural network.
\end{eabstract}
